\documentclass[10pt,a4paper]{amsart}
\usepackage[margin=19mm]{geometry}
\usepackage{amsmath,amssymb,mathtools}
\usepackage[scr=rsfs,cal=euler]{mathalfa}
\usepackage{xfrac}
\usepackage[unicode,hidelinks,bookmarks=false]{hyperref}
\newcommand{\ie}{i.e.\ }
\newcommand{\cf}{cf.\ }
\newcommand{\resp}{respectively\ }
\newcommand{\Subsection}{\subsection}
\providecommand{\nobreakdash}{\nobreak\hbox{-}}
\setlength{\emergencystretch}{2em}
\multlinegap0pt
\usepackage{amssymb}
\usepackage{enumerate}
\usepackage{tikz}
\usepackage[all]{xy}
\usepackage{xcolor}
\newtheorem{theorem}{Theorem}
\title{When is the canonical conductor minimal?}
\author{\"Ozg\"ur Esentepe}
\date{Source: arXiv:2509.22966v2 (CC BY 4.0)}
\begin{document}
\maketitle

\noindent\emph{Source and licence.} When is the canonical conductor minimal?. Version arXiv:2509.22966v2 (CC BY 4.0), distributed by arXiv under the Creative Commons Attribution 4.0 International licence, http://creativecommons.org/licenses/by/4.0/, as recorded in the arXiv licence metadata for this exact version and shown on the abstract page https://arxiv.org/abs/2509.22966v2. Full licence notice: \texttt{licenses/arxiv-2509.22966-CC-BY-4.0.txt}.\par\smallskip

\noindent\textit{Editorial scope.} Counted unit: the article's theorem pseudo-frobenius-theorem (source0.tex lines 426--428). The article's complete original proof of it is delivered with it (source0.tex lines 429--437, one \texttt{proof} environment of 9 lines). No mathematics is written, repaired or re-derived: the statement and the proof are the article's own lines, in the article's own notation, and no retained line is substituted or re-flowed.  No further source-local context is needed: every cross-reference of every retained line resolves inside the two delivered spans.  Nothing was omitted from any retained span, so the package carries no omission record.  No retained line contains a citation, so the wrapper carries no bibliography and no bibliography entry.  No retained line contains a figure environment, an image-inclusion command or drawing code, so no image is delivered and the package declares no asset dependency.  Editorial formatting only, in addition: an AMS \texttt{amsart} preamble that restates the article's own declarations and macros used by the retained lines, namely the environments \texttt{theorem} on separate counters, together with the standard packages those lines need.  The retained environments print in this document as Theorem 1 in order of appearance, whereas the article prints the same items with the numbering of its own full document.  The complete native source of the article as distributed in the arXiv e-print for this exact version is bundled unchanged as \texttt{sources/185852/source0.tex}.
\begin{theorem}\label{pseudo-frobenius-theorem}
    Let $H$ be a numerical semigroup with Frobenius number $f$. Then, the ring $R = k[[H]]$ has minimal canonical conductor if and only if $f-1$ is a pseudo-Frobenius number for $H$. 
\end{theorem}

\begin{proof}
    If $f-1$ is a pseudo-Frobenius number for $H$, then there exists a canonical fractional ideal $\omega = R + Rt + J$ where $J = R^{t_1} + \ldots + R^{t_s}$ for some $t_1, \ldots, t_s$. Note that $\omega^n \subseteq k[[t]] = \overline{R}$ for any positive integer $n$. On the other hand, we have 
    \begin{align*}
    (R+R^t)^n = R + Rt + Rt^2 + \cdots + Rt^n \subseteq \omega^n.
    \end{align*}
    For large enough $n$ (we can choose $n \leq f$), we have $(R+Rt)^n = k[[t]] = \overline{R}$. Hence, we obtain an isomorphism $\omega^n \cong \overline{R}$.

    If $f-1$ is not a pseudo-Frobenius number, we the canonical fractional ideal is of the form $\omega = R + R^{t_1} + \cdots + R^{t_s}$ for some $t_1, \ldots, t_s \geq 2$. Therefore, we have $t \notin \omega^n$ for any $n \geq 1$. This finishes the proof. 
\end{proof}

\end{document}
